\documentclass[12pt]{article}
\usepackage[utf8]{inputenc}
\usepackage[T1]{fontenc}
\usepackage[margin=1in]{geometry}
\usepackage{setspace}
\usepackage{xcolor}
\usepackage[colorlinks=true, urlcolor=blue!60!black, linkcolor=blue!60!black]{hyperref}
\setlength{\parindent}{0em}
\setlength{\parskip}{0.45em}
\renewcommand{\baselinestretch}{1.0}
\begin{document}
\begin{center}
{\Large \textit{Moore v. United States}, 602 U.S. 572 (2024)\par}
\vspace{0.3em}
{\normalsize Case Brief --- One-Page Edition\par}
\end{center}

\textbf{Memory Jogger.} Congress may tax shareholders on an entity's realized, undistributed income without apportionment (\href{https://supreme.justia.com/cases/federal/us/602/22-800/}{Justia}).

\textbf{Facts.} The Moores owned 13\% of KisanKraft, an Indian CFC that reinvested its profits. The 2017 Mandatory Repatriation Tax taxed them on their share of its accumulated earnings.

\textbf{Procedural History.} Their refund suit was dismissed; the Ninth Circuit affirmed, saying ``realization of income is not a constitutional requirement.''

\textbf{Votes.} Affirmed, 7--2 (Kavanaugh for the Court; Jackson concurring; Barrett and Alito in the judgment; Thomas and Gorsuch dissenting).

\textbf{Holding.} The MRT taxes income realized by the CFC and attributed to its shareholders, which Congress may do without apportionment when the entity is untaxed; wealth taxes and unrealized appreciation were reserved.

\textbf{Reasoning.} The Sixteenth Amendment restored income taxes to the unapportioned class, and precedents and practice since 1864 allow taxing owners on entity income; the Moores' theory would cause ``fiscal calamity.'' \emph{Separately:} Jackson saw no realization requirement; Barrett found one but concurred because the Moores conceded Subpart F; Thomas would require realization by the taxpayer.

\textbf{Examples.} \emph{Same-Side:} partners taxed on retained profits. \emph{Opposite-Side:} a tax on billionaires' unrealized gains. \emph{Fence-Sitter:} pass-through taxation of public shareholders.

\textbf{Critique.} Faulted for ducking the question presented; formalism versus economic substance.

\textbf{Key Quotation.} ``So the precise and narrow question that the Court addresses today is whether Congress may attribute an entity's realized and undistributed income to the entity's shareholders or partners, and then tax the shareholders or partners on their portions of that income.'' --- Justice Kavanaugh, Opinion of the Court.

\end{document}
